\section{Mixes and Medleys}

\begin{question}[subtitle = {January 2026 Problem 1, edited},ID=Jan26_1]
On this problem, only your answer will be graded. No justification is required.
\begin{enumerate}
    \item How many ring homomorphisms
    \[\Z[x,x\inv] \to \F_{27}\]
    are there? Here, $\Z[x,x\inv]$ is the localization of $\Z[x]$ inverting $x$, and $\F_{27}$ is the finite field with 27 elements.
    \item Given $n \geq 2$, how many group homomorphisms
    \[S_n \to \Z/2\Z\]
    are there?
    \item $V$ is an $n$-dimensional vector space over a field $K$. Let $W$ be the vector space of alternating $(n+1)$-linear forms $V \times V \times \dots \times V \to K$ (in $n+1$ vector arguments). What is $\dim_K(W)$?
    \item Suppose $K \supset E \supset F$ is a tower of field extensions such that both $K \supset E$ and $E \supset F$ are finite Galois extensions. Is it true that $K$ must be a finite Galois extension of $F$? %original problem said ``of E''
    \item Find the cardinality of the tensor product
    \[((\Z/20\Z)\oplus (\Z/26\Z)) \tensor_{\Z} ((\Z/10\Z)^2).\]
\end{enumerate}
\end{question}
\begin{solution}
\begin{enumerate}
    \item A ring homomorphism must send 1 to 1, so a homomorphism $\Z[x,x\inv] \to \F_{27}$ is completely determined by the image of $x$. We can send $x$ to any invertible element of $\F_{27}$, of which there are 26 choices, so there are 26 different homomorphisms. (The biggest danger here is inadvertantly mentally replacing $\F_{27}$ with $\Z/27\Z$, ask me how I know.)
    \item Two: the zero homomorphism and the sign homomorphism (the quotient $S_n \to S_n/A_n$). Indeed, suppose $\phi\colon S_n \to \Z/2\Z$ is nontrivial, hence surjective. For any $\sigma \in S_n$, we have $\phi(\pi \sigma \pi\inv) = \phi(\sigma)$ for any $\pi \in S_n$, because $\Z/2\Z$ is commutative. If $\tau_1, \tau_2 \in S_n$ are transpositions, then they are conjugate, and hence $\phi(\tau_1) = \phi(\tau_2)$. Since $S_n$ is generated by transpositions, we have that $\phi(\tau)$ generates the $\im \phi$ for any transposition $\tau$. We assumed $\phi$ is surjective, so $\phi(\tau) = 1$. Thus, all even permutations are in $\ker \phi$, i.e.\ $\ker \phi \supset A_n$. Since $A_n$ is already index 2, we must have $\ker \phi = A_n$.
    \item This is the first time the language of alternating multilinear forms has appeared on the qual, so I'll include some definitions. An $m$-linear form is a function $f$ which takes $m$ vector inputs $v_1,\dots,v_m \in V$ and outputs an element of $K$ such that $f$ is linear in each argument separately: $f(v_1,\dots,av_i + b v_i',\dots,v_m) = a f(v_1,\dots,v_i,\dots,v_m) + b f(v_1,\dots,v_i',\dots,v_m)$ for each index $1 \leq i \leq m$. An $m$-linear form is called ``alternating'' if it evaluates to 0 whenever two or more of its arguments are the same, e.g.\ if $m=2$ an alternating form satisfies $f(v,v)=0$ for all vectors $v$. The determinant is an example of an alternating $n$-linear form. The set of all $m$-linear forms on $V$ is a vector space over $K$, and the set of alternating forms is a subspace.
    
    The key fact, which is an exercise if you don't already know it, is that an alternating form evaluates to 0 whenever its arguments are linearly dependent. Since $V$ is dimension $n$, any set of $n+1$ vectors will be dependent, and thus any alternating $(n+1)$-linear form will evaluate to 0 on all inputs. Thus, $\dim_K(W) = 0$.
    \item $K$ need not be a Galois extension of $F$ (but will be finite). As written, the problem didn't say to give a counterexample, but you should anyway. The standard counterexample is $\Q(\sqrt[4]{2}) \supset \Q(\sqrt{2}) \supset \Q$.
    \item We use two tensor product facts: tensor products distribute over direct sums (justifying the notation), and $(\Z/a\Z) \tensor_{\Z} (\Z/b\Z) \cong \Z/\gcd(a,b)\Z$. Thus:
    \begin{align*}
        ((\Z/20\Z)\oplus (\Z/26\Z)) \tensor_{\Z} ((\Z/10\Z)^2) &\cong ((\Z/20\Z) \tensor_{\Z} (\Z/10\Z)^2) \oplus ((\Z/26\Z) \tensor_{\Z} (\Z/10\Z)^2)\\
        &\cong ((\Z/20\Z) \tensor_{\Z} (\Z/10\Z))^2 \oplus ((\Z/26\Z) \tensor_{\Z} (\Z/10\Z))^2\\
        &\cong (\Z/10\Z)^2 \oplus (\Z/2\Z)^2,
    \end{align*}
    which has order $10^2 \cdot 2^2 = 400$.
\end{enumerate}
\end{solution}

% TENS, LOC
\begin{question}[subtitle = August 2025 Problem 1,ID=Aug25_1]
On this problem, only your answer will be graded. No justification is required.
\begin{enumerate}
    \item Suppose $x_1,\dots,x_n$ are $n$ distinct complex numbers, and $a_1,\dots,a_n \geq 0$, $b_1,\dots,b_n \geq 0$ are integers. Put
    \[f(x) = (x-x_1)^{a_1} \cdots (x-x_n)^{a_n} \qquad g(x) = (x-x_1)^{b_1} \cdots (x - x_n)^{b_n}.\]
    (Some of the exponents $a_i, b_i$ may be zero, so $f$ and $g$ need not share all roots.)
    Consider $\C[x]$-modules $M = \C[x]/(f(x))$ and $N = \C[x]/(g(x))$. Write a formula for
    \[\dim_{\C} (M \tensor_{\C[x]} N)\]
    in terms of $a_1,\dots,a_n$ and $b_1, \dots, b_n$.
    \item In the situation of question (a), write a formula for $\dim_{\C} M_g$ in terms of $a_1,\dots,a_n$ and $b_1,\dots,b_n$. Here $M_g$ is the localization of $M$ with respect to $g \in \C[x]$.
    \item Let $M$ be the $\C[x]$-module
    \[\C[x]/(x^3-x) \times \C[x]/(x^3+x^2) \times \C[x]/(x^2+1).\]
    Rewrite $M$ in the form
    \[M \simeq \C[x]/(p_1(x)) \times \C[x]/(p_2(x)) \times \cdots \times \C[x]/(p_k(x))\]
    for non-constant polynomials $p_1(x),\dots,p_k(x) \in \C[x]$ such that $p_k(x) \mid p_{k-1}(x) \mid \dots \mid p_1(x)$ (that is, the $p_i(x)$ are the \emph{invariant factors} of $M$).
    \item In the situation of question (c), consider $M$ as a vector space over $\C$ via restriction of scalars to $\C \subset \C[x]$. Let the operator $T\colon M \to M$ be the action of $x \in \C[x]$ on $M$. Clearly, it is a $\C$-linear operator; what are its eigenvalues?
    \item Give an example of a ring $R$, an $R$-module $N$, and an exact sequence of $R$-modules
    \[0 \to M_1 \to M_2 \to M_3 \to 0\]
    such that the induced sequence
    \[0 \to \Hom_R(N,M_1) \to \Hom_R(N,M_2) \to \Hom_R(N,M_3) \to 0\]
    is \underline{not} exact.
\end{enumerate}
\end{question}
\begin{solution}
This is a reading-intensive problem 1.
\begin{enumerate}
    \item By standard tensor product properties, $R/I \tensor_R R/J \cong R/(I+J)$. In a PID like $\C[x]$, the sum of two ideals is just the ideal generated by the \emph{gcd} of the two generators: $(f(x)) + (g(x)) = (\gcd(f(x),g(x))$. From the factorizations given, we can see that
    \[\gcd(f(x),g(x)) = (x-x_1)^{\min(a_1,b_1)} \cdots (x-x_n)^{\min(a_n,b_n)}.\]
    Finally, the dimension of $\C[x]/(\mathrm{tomato}(x))$ is just $\deg(\mathrm{tomato}(x))$. Thus, $\dim_{\C}(M \tensor_{\C[x]} N) = \sum_{i=1}^n \min(a_i,b_i)$.
    \item (This is a bit like \citeQ{Jan23_1}(a).) Localizing with respect to $g$ is the same as localizing with respect to every factor $(x-x_i)$ satisfying $b_i > 0$. Thus,
    \begin{align*}
    M_g &= \paren{\C[x]/(f(x))}_g\\
    &\cong \C[x]_g/(f(x))_g\\
    &\cong \C[x]_g/(h(x))_g,
    \end{align*}
    where $h(x)$ has the factorization of $f(x)$ without any of the factors $(x-x_i)^{a_i}$ that have $b_i > 0$. So, $\displaystyle \dim_{\C}(M_g) = \sum_{i: b_i = 0} a_i$.
    \item We start by factoring these polynomials: $x^3-x = x(x-1)(x+1)$, $x^3+x^2 = x^2(x+1)$, $x^2+1 = (x-i)(x+i)$. For each irreducible polynomial appearing in these factorizations, record the powers to which it appears in descending order, as shown in the table below.

    \begin{minipage}{.2\textwidth}
    \begin{tabular}{c|l}
        poly. & powers \\
        \hline
        $x$ & 2, 1 \\
        $x-1$ & 1 \\
        $x+1$ & 1, 1 \\
        $x-i$ & 1 \\
        $x+i$ & 1
    \end{tabular}
    \end{minipage}
    \begin{minipage}{.75\textwidth}
    To find the polynomial $p_i(x)$, we go through our table and include a factor of $(x-a)^{e_i}$, where $e_i$ is the $i$th power appearing on the right side of the table (or 0 if the list is too short). So, we have $p_1(x) = x^2(x-1)(x+1)(x-i)(x+i) = x^6-x^2$ and $p_2(x) = x(x+1) = x^2+x$, and then we are done since none of the lists of exponents have more than 2 entries. We can verify that $p_2(x) \mid p_1(x)$, and use the Chinese Remainder Theorem to verify that $M = \C[x]/(x(x+1)) \times \C[x]/(x^2(x-1)(x+1)(x-i)(x+i))$.
    \end{minipage}
    \item The minimal polynomial of $T$ is $p_1(x) = x^2(x-1)(x+1)(x-i)(x+i)$, so the eigenvalues of $T$ are $0,\pm 1, \pm i$.
    \item Take $N = \Z/2\Z$, and the short exact sequence
    \[0 \to \Z \tonamed{\times 2} \Z \to \Z/2\Z \to 0.\]
\end{enumerate}
\end{solution}

% Part (c) originally said "what is the dimension" without specifying what field they were thinking of it as being over. I decided the problem is slightly more interesting if you say you want the real dimension.
% TENS
\begin{question}[subtitle = {January 2025 Problem 1, edited},ID=Jan25_1]
On this problem, no justification is required; only the answers will be graded.
\begin{enumerate}
    \item Let $R$ be a (not necessarily commutative) ring, and let $M$ be a right $R$-module. Let $A$ be the annihilator of $M$ in $R$. Is $A$ necessarily a two-sided ideal of $R$, necessarily a one-sided ideal of $R$, or not necessarily either of these things?
    \item Let $H$ be the subgroup of $S_5$ generated by the element $(12)(345)$. Then $H$ acts on $S_5$ by left multiplication. How many orbits are there of this action?
    \item Let $\mbb{H}$ be the quaternion algebra $\R + \R i + \R j + \R k$, and let $A$ be its subalgebra $\R + \R i$ (so $A$ is isomorphic to the complex numbers). What is the dimension of $\mbb{H} \tensor_A \mbb{H}$ as a real vector space?
    \item True or false: if $A$ is an $n \times n$ orthogonal matrix, then all $n$ eigenvalues are either $1$ or $-1$.
    \item Give an example of a prime ideal in the ring $\C[x,y,z]/(xy,xz,yz,x^2,z^2)$ which is not a maximal ideal.
\end{enumerate}
\end{question}
\begin{solution}
\begin{enumerate}
    \item $A$ is always a two-sided ideal of $R$. Indeed, for any $r \in R$, $a \in A$, and $m \in M$, $(ra)m = r(am) = r\cdot 0 = 0$, so $ra \in A$. Also, $(ar)m = a(rm) = 0$ since $rm \in M$ and $a$ annihilates $M$, so $ar \in A$. It is easy to see that $A$ is closed under addition.
    \item Orbits are cosets of $H$, so there are $[S_5:H] = 20$ orbits.
    \item 8
    \item False, e.g.\ $\begin{pmatrix*}[r] 0 & 1 \\ \shortminus 1 & 0 \end{pmatrix*}$
    \item Prime ideals of this ring correspond to prime ideals of $\C[x,y,z]$ which contain $(xy,xz,yz,x^2,z^2)$. Notice that every generator of that ideal is divisible by $x$ or by $z$, so $(x,z)$ is a prime ideal containing that ideal. On the other hand, $(x,z)$ is not maximal because e.g.\ $(x,y,z)$ is another prime ideal which is strictly bigger.
\end{enumerate}
\end{solution}

%
\begin{question}[subtitle=August 2024 Problem 1,ID=Aug24_1]
On this problem, no justification is required; only the answers will be graded.
\begin{enumerate}
    \item The symmetric group on 5 letters $S_5$ acts on the set $\set{1,\dots,5} \times \set{1,\dots,5} \times \set{1,\dots,5}$ by the diagonal action $\sigma(a,b,c) = (\sigma(a),\sigma(b),\sigma(c))$. How many orbits are there of this action?
    \item The vector space $\R^3$ is a module for the ring of matrices $M_3(\R)$. Let $v$ be a nonzero vector in $\R^3$. What is the dimension of the annihilator of $v$ in $M_3(\R)$?
    \item Let $\F_5$ be the field with 5 elements. How many ideals does the ring $\F_5[T]/(T^2+1)$ have?
    \item Give an example of an abelian subgroup of $\GL_2(\C)$ that is not contained in the group of diagonal matrices. %Like Jan20_5 (b) but the opposite
    \item Give an example of a simple ring which is not a division ring. %repeat from January
\end{enumerate}
\end{question}
\begin{solution}
\begin{enumerate}
    \item There are 5 orbits: 1 orbit where none of the three coordinates are equal, 3 orbits where two of the three coordinates are equal and the third is distinct, and 1 orbit where all three coordinates are equal.
    \item If we extend the set $\set{v}$ to a basis $\set{v,v_2,v_3}$, then the annihilator of $v$ is the same as the set of matrices whose first column consists entirely of 0s. This is a subspace of dimension 6 inside $M_3(\R)$.
    \item By the lattice isomorphism theorem, ideals in $\F_5[T]/(T^2+1)$ correspond to ideals in $\F_5[T]$ that contain $(T^2+1)$. Since $\F_5[T]$ is a UFD, this corresponds to the number of divisors of the polynomial $T^2+1$ over $\F_5$. Over $\F_5$, the polynomial $T^2+1$ factors as $(T-2)(T-3)$, so it has four total divisors, giving us the four containing ideals: $(1), (T-2), (T-3), (T^2+1)$.
    \item The subgroup generated by
    \[\twomatrix{1}{1}{0}{1}\]
    forms an abelian group isomorphic to $\Z$, but which is obviously not contained in the group of diagonal matrices.
    \item Any counterexample must be a noncommutative ring, because a simple commutative ring must be a field. The standard example is $M_2(\R)$, which is simple because its only two-sided ideals are the zero ideal and the whole ring, but it is not a division ring because there are matrices that don't have inverses.
\end{enumerate}
\end{solution}

%a mean one!
%
\begin{question}[subtitle=January 2024 Problem 1,ID=Jan24_1]
On this problem, no justification is required; only the answers will be graded.
\begin{enumerate}
    \item Give an example of a simple ring which is not a division ring.
    \item Let $A \colon \R^3 \to \R^3$ be the linear transformation defined by $Ae_1 = e_1 + e_2$, $Ae_2 = e_2$, $Ae_3 = e_3$. Let $V$ be the vector space of linear transformations from $\R^3$ to $\R^3$ which commute with $A$. What is $\dim_{\R}(V)$?
    \item Give an example of a group $G$ and a subgroup $H$, and a subgroup $K$ of $H$, such that $K$ is normal in $H$, and $H$ is normal in $G$, but $K$ is not normal in $G$.
    \item Give an example of an inseparable field extension.
    \item Give an example of a pair $(R,M)$, where $R$ is a ring and $M$ is an $R$-module that is not Noetherian.
\end{enumerate}
\end{question}
\begin{solution}
\begin{enumerate}
    \item (deja vu) Any counterexample must be a noncommutative ring, because a simple commutative ring must be a field. The standard example is $M_2(\R)$, which is simple because its only two-sided ideals are the zero ideal and the whole ring, but it is not a division ring because there are matrices that don't have inverses.
    \item $\dim_{\R}(V) = 5$. One way to do this problem is to write out the matrix of $A$, and an arbitrary matrix $B$, and compute the linear relations that must hold in order for them to commute. On the qual, if I saw a problem like this and didn't think of a clever workaround immediately, I would try this first.
    
    Here is a trick that comes up often with commuting matrices which does not outright solve this problem, but does reduce the amount of computation required: if $v$ is an eigenvector of $A$ of eigenvalue $\lambda$, and $B$ commutes with $A$, then $Bv$ is \emph{also} an eigenvalue of $A$ of eigenvalue $\lambda$. Indeed, $ABv = BAv = \lambda Bv$. This immediately tells us that for $i = 2,3$, $Be_i = b_{2i}e_2 + b_{3i} e_3$, i.e. that the latter two entries in the first row of $B$ are 0. To finish this problem, I tried doing a similar computation for $e_1$: $ABe_1 = BAe_1 = Be_1 + Be_2$, which we can rearrange to get $(A-I)Be_1 = Be_2$. Here, writing out the matrix $A-I$ is very easy, and I quickly saw that $b_{11} = b_{22}$ and $b_{32} = 0$ were the last two linear conditions defining $B$.
    \item Any counterexample must be a nonabelian group. The two standard examples are $D_4 = \genset{r,s: r^4=s^2=e, srs=r\inv}$, with $H = \genset{r^2,s} \cong \Z/2\Z \times \Z/2\Z$ and $K = \genset{s}$, and $A_4$, with $H = \genset{(12)(34),(13)(24)} \cong \Z/2\Z \times \Z/2\Z$ and $K = \genset{(12)(34)}$. In fact, for any even number $n = 2m$, $D_n$ is a counterexample with $H = \genset{r^m,s}$ and $K = \genset{s}$. (If you're curious about cases where transitivity holds, look for ``characteristic subgroup''.)
    \item Any example must be over a non-perfect field, so in particular must be over an infinite field of characteristic $p \neq 0$. The standard example is to take the base field to be the field of rational function $\F_p(t)$, and the extension to be $\F_p(\sqrt[p]{t}) \cong \F_p(t)[x]/(x^p-t)$.
    \item Noetherian modules are less familiar than Noetherian rings are: a module $M$ is Noetherian if its submodules satisfy the ACC, which is equivalent to saying every submodule is finitely generated. Maybe the easiest example of a non-Noetherian module would be to take $R$ to be a non-Noetherian ring, and $M$ to be a non-finitely-generated ideal in $R$. For example, if $R = \C[x_1,x_2,x_3,\dots]$ is the polynomial ring in countably many variables, and $M = (x_1,x_2,x_3,\dots)$ is the (maximal) ideal generated by the individual variables. If you don't like that, you can also take any infinitely-generated module over your favorite ring of choice, e.g. $R = \Z$, $M = \Z^{\oplus \N}$.
\end{enumerate}
\end{solution}

% TENS, LOC
\begin{question}[subtitle=August 2023 Problem 1,ID=Aug23_1]
On this problem, no justification is required; only the answers will be graded.
\begin{enumerate}
    \item Give an example of a field $K$, a finite algebraic field extension $L/K$, and a finite-dimensional $K$-algebra $A$ such that $A$ has no zero divisors but $A \tensor_K L$ does have zero divisors.
    \item Give an example of a zero divisor in the $A \tensor_K L$ you provided above.
    \item Give two elements $x,y$ of the algebra of quaternions over $\R$ such that $x$ and $y$ both have norm 4 and $x$ and $y$ do not commute.
    \item Give an example of a ring with exactly three ideals (including the zero ideal and the ring itself). %soft...
    \item Give an example of a localization of $\Z$ which is not of the form $\Z[1/n]$ for any integer $n$.
\end{enumerate}
\end{question}
\begin{solution}
\begin{enumerate}
    \item The example that came to mind for me is that if $L/K$ is a finite algebraic field extension, say $L = K(\alpha)$, then $L$ is moreover a finite-dimensional $K$-algebra. If we tensor $L \tensor_K L$ we will get $L \tensor_K K[x]/m_{\alpha}(x) \cong L[x]/m_{\alpha}(x)$, where $m_{\alpha}(x)$ is the minimal polynomial of $\alpha$. As a concrete example, let $K = \R$, $L = \C$, then $\C \tensor_{\R} \C \cong \C \tensor_{\R} \R[x]/(x^2+1) \cong \C[x]/(x^2+1)$, which has zero-divisors because in this ring $(x-i)(x+i) = 0$.
    \item See above.
    \item We could take $x = 2i$ and $y = 2j$.
    \item You have lots of options for how to come up with examples here, see \citeQ{Jan20_3}.
    \item I think this is a confusingly-phrased question, it took me a while to figure out what it was asking. What we want is a multiplicatively closed subset $S \subset \Z$ so that $S\inv \Z$ is not of the form $\Z[1/n]$ for some integer $n$. One idea might be to try throwing in two different denominators, say $\Z[1/2,1/3]$, but this doesn't work because $\Z[1/2,1/3] = \Z[1/6]$. For the same reason, throwing in only finitely many denominators doesn't work. What would work would be something like $\Z_{(2)}$, which as a reminder is where we throw in all denominators of numbers not in the ideal $(2)$. Certainly, there is no way for us to get every odd denominator by only taking a ring of the form $\Z[1/n]$.
\end{enumerate}
\end{solution}

% TENS, LOC
\begin{question}[subtitle=January 2023 Problem 1,ID=Jan23_1]
On this problem, only your answer will be graded, no justification is necessary.
\begin{enumerate}
    \item Let $R$ be a commutative ring and take $f,g \in R$. Consider the homomorphism \mbox{$R \to R/(f) \tensor_R R_g$.} Describe its kernel in terms of $f$ and $g$.
    \item Consider the homomorphism
    \[\phi\colon (\Z/125\Z) \times (\Z/25\Z) \to \Z/5\Z\colon (a,b) \mapsto a+b.\]
    (Note that $a+b$ makes sense as an element of $\Z/5\Z$.) The kernel $\ker \phi$ is a finite abelian group. What is its isomorphism class. (That is, its canonical form, such as $(\Z/125\Z) \times (\Z/25\Z)$.)
    \item True or False: if $\alpha$ and $\beta$ are algebraic numbers such that $\Q(\alpha)$ is a Galois extension of $\Q$, and $\Q(\alpha,\beta)$ is a Galois extension of $\Q(\alpha)$, then $\Q(\alpha,\beta)$ is a Galois extension of $\Q$.
    \item True or False: For any non-zero $f \in \C[x]$, any finitely-generated torsion-free module over $\C[x,\frac{1}{f}]$ is free.
    \item True or False: For any $n$, any ideal $I \subset \C[x_1,\dots,x_n]$ can be generated by (at most) $n$ elements.
\end{enumerate}
\end{question}
\begin{solution}
\begin{enumerate}
    \item It's important here that localization is exact, so $R/(f) \tensor_R R_g \cong R_g/(f)_g$, and the homomorphism we're looking at is just the composition of $R \to R_g \to R_g/(f)_g$. Certainly, we kill all the multiples of $f$, but because we're allowed to divide by $g$ in $R_g$, the kernel will also contain any element that can be ``multiplied into'' the ideal $(f)$ by some power of $g$. As a concrete example, suppose $R = \Z$, $f = 6$, $g = 9$, then 2 will be in the kernel of $\Z \to \Z[1/9]/6 \Z[1/9]$ because $2 = 18/9 = 6(3/9) \in 6\Z[1/9]$.
    
    (I think writing something like the above qualifies as a complete answer already, but I'll remark that the set in question $\set{r \in R: \exists n,\ g^n r \in (f)}$ has a name: it's the \emph{ideal saturation} of the ideal $(f)$ by the ideal $(g)$, denoted $(f): (g)^{\infty}$.)
    \item We can see that the group on the left has order $5^5$, and the group on the right has order 5, and the map is certainly surjective, so the kernel must have order $5^4$, and is a subgroup of $(\Z/125\Z) \times (\Z/25\Z)$, so the possibilities are either $(\Z/125\Z) \times (\Z/5\Z)$ or $(\Z/25\Z) \times (\Z/25\Z)$. We can check which it is by checking to see if the kernel has an element of order 125, which in fact it does, for example $(1,4)$.
    \item False. The standard counterexample is $\alpha = \sqrt{2}$, $\beta = \sqrt[4]{2}$.
    \item True. $\C[x]$ is a PID, hence all its localizations are also PIDs, and a finitely generated modules over a PID is torsion-free iff it is free.
    \item False. The standard example is $(x^2,xy,y^2) \in \C[x,y]$.
\end{enumerate}
\end{solution}

%LGFF
% Commuting-ish matrices
\begin{question}[subtitle = January 2023 Problem 5,ID=Jan23_5]
Let $p$ be a prime number and let $G$ be the group of $2 \times 2$ matrices of the form
\[\begin{pmatrix}1&a\\0&b\end{pmatrix},\qquad a \in \F_p,\ b\in \F_p\mult.\]
Fix a dimension $n$, and let $\rho\colon G \to \GL_n(\C)$ be a homomorphism. Put
\[x = \begin{pmatrix}1&1\\0&1\end{pmatrix},\ y_b = \begin{pmatrix}1&0\\0&b\end{pmatrix}\qquad (b \in \F_p\mult).\]
\begin{enumerate}
    \item Show that $\rho(x)$ is diagonalizable, and that its eigenvalues are $p$th roots of unity.
    \item Suppose $v \in \C^n$, $v \neq 0$ is an eigenvector of $\rho(x)$ of eigenvalue $\lambda$. Prove that for any $b \in \F_p\mult$, $\rho(y_b)v$ is also an eigenvector of $\rho(x)$, and find its eigenvalue.
    \item Let $v$ be as in part (b). Show that $W = \Span\genset{\rho(y_b)v: b \in \F_p\mult}$ is $G$-invariant (that is, $\rho(g)W \subseteq W$ for all $g \in G$), and that, if $\lambda \neq 1$, $\dim W = p-1$.
\end{enumerate}
\end{question}
\begin{solution}
\begin{enumerate}
    \item Note that in $G$, $x^p = I$, so $\rho(x)$ satisfies the polynomial $x^p-1$. Thus, all the eigenvalues of $\rho(x)$ must be roots of $x^p-1$, i.e.\ $p$th roots of unity. Furthermore, the minimal polynomial of $\rho(x)$ divides $x^p-1$, which has all distinct roots, so the minimal polynomial must have all distinct roots, so $\rho(x)$ must be diagonalizable.
    \item In $G$, we have that $y_b x^b = x y_b$. Thus, $\rho(x) \rho(y_b) v = \rho(y_b) \rho(x)^b v = \lambda^b \rho(y_b) v$. Thus, $\rho(y_b) v$ is an eigenvector of eigenvalue $\lambda^b$.
    \item First note that $x$ and the set of all $y_b$ together generate $G$, since
    \[\begin{pmatrix}1&a\\0&b\end{pmatrix} = y_b x^a.\]
    Thus, since each of the generators of $W$ is an eigenvector of $\rho(x)$, and these eigenvectors are permuted by each $\rho(y_b)$, the generating set of $W$ is fixed by the action of $G$, and hence $W$ is $G$-invariant. If $\lambda \neq 1$, $\lambda$ is a primitive $p$th root of unity, and so each $\lambda^b$ is distinct. Thus, the $\rho(y_b)v$ are linearly independent, since they are eigenvectors of $\rho(x)$ of different eigenvalues. Since there are $p-1$ choices for $b$, this means $\dim W = p-1$.
\end{enumerate}
\end{solution}

%non-Galois extension, definition of radical of an ideal, a group of order p^2 is abelian + structure theorem, tensor product + torsion, simplicity of M_2(\C)
% TENS
\begin{question}[subtitle=August 2022 Problem 1,ID=Aug22_1]
On this problem, only your answer will be graded, no justification is necessary.
\begin{enumerate}
    \item Let $K = \Q(2^{1/3})$. Is $K/\Q$ a Galois extension?
    \item What is the radical of the ideal $(x^3-x^2) \subset \R[x]$?
    \item How many different (up to isomorphism) groups of order 49 are there?
    \item Give an example of two non-zero abelian groups $M, N$ such that $M \tensor_\Z N = 0$.
    \item How many two-sided ideals are there in $M_2(\C)$, the ring of $2 \times 2$ matrices with entries in $\C$?
\end{enumerate}
\end{question}
\begin{solution}
\begin{enumerate}
    \item No, since $\zeta_3 2^{1/3}$ is a root of $x^3-2$ which does not lie in $K$.
    \item $(x^2-x)$, in this case radical corresponds to ``largest squarefree divisor''.
    \item 2, $\Z/49\Z$ and $\Z/7\Z \times \Z/7\Z$.
    \item $M = \Z/2\Z$, $N = \Z/3\Z$.
    \item 2, the zero ideal and the whole ring.
\end{enumerate}
\end{solution}

%Jan '22 was all about S_4

%LGFF. This problem involves Jordan NF in part (a), and field theory for parts (c) and (d).
\begin{question}[subtitle = January 2022 Problem 2,ID=Jan22_2]
For each of the following statements, either prove it or provide a counterexample.
\begin{enumerate}
    \item Every matrix of finite order in $\GL_n(\C)$ is diagonalizable.
    \item Every matrix of finite order in $\GL_n(\Q)$ is diagonalizable.
    \item Every matrix in $\GL_n(\widebar{\F}_p)$ has finite order. (We write $\widebar{\F}_p$ for the algebraic closure of the finite field $\F_p$.)
    \item Every matrix of finite order in $\GL_n(\widebar{\F}_p)$ is diagonalizable.
\end{enumerate}
\end{question}
\begin{solution}
\begin{enumerate}
    \item Suppose $A \in \GL_n(\C)$ has finite order, so that $A^k = I$ for some $k$. Then we know the minimal polynomial of $A$ divides $x^k - 1$. Now, over $\C$, $x^k-1$ splits into $k$ different linear factors, so the minimal polynomial also splits into different linear factors. Since the exponents on the linear factors of the minimal polynomial are the sizes of maximal Jordan blocks, this tells us that the size of a maximal Jordan block is 1, which is the same as saying $A$ is diagonalizable.
    \item Everyone's favorite counterexample:
    \[\begin{pmatrix}0 & \shortminus 1\\ 1 & 0\end{pmatrix}.\]
    \item Let $\mu_A(x)$ be the minimal polynomial of $A$. Let $\F_q \supset \F_p$ be the smallest extension containing all the roots of $\mu_A(x)$. Then all the roots of $\mu_A(x)$ are nonzero elements of $\F_q$, and so satisfy $\alpha^{q-1} = 1$. Since $\mu_A(x)$ factors over $\F_q$, write \[\mu_A(x) = (x-\alpha_1)^{e_1} \cdots (x-\alpha_k)^{e_k}.\]
    If all the $e_i$ were 1 then we would just know that $\mu_A(x) \mid x^{q-1}-1$, and we would be done, but if one of the $e_i > 1$ then we want to find another polynomial of the form $x^{\text{something}} - 1$ which $\mu_A$ divides. We know $x^{q-1} - 1$ has distinct roots, so if $e = \max_i(e_i)$, then $\mu_A(x) \mid (x^{q-1}-1)^e$, but this isn't in the form we want yet. To get it in the desired form, we can instead choose the smallest power of $p$ which is greater than $e$, since powers of $p$ distribute over addition in characteristic $p$. Namely, set $m = \ceil{\log_p(e)}$, then we have $\mu_A(x) \mid (x^{q-1}-1)^{p^m} = x^{p^m(q-1)} - 1$. Thus, $A^{p^m(q-1)} = I$. (I actually think this exponent might be optimal, but I'm not sure. Lemme know if you figure out one way or the other.)
    \item The problem with our proof from part (a) is that we said $x^k-1$ splits into \emph{distinct} linear terms over $\C$, but that need not be true in $\widebar{\F}_p$. For example, $x^p-1$ definitely doesn't. So, take the example of
    \[\begin{pmatrix} 1 & 1\\0 & 1 \end{pmatrix}\]
    over $\widebar{\F}_2$.
\end{enumerate}
\end{solution}

%2 questions about localization, 1 somewhat weird Jordan form computation
% LOC
\begin{question}[subtitle = August 2021 Problem 1,ID=Aug21_1]
On this problem, only the answer will be graded.
\begin{enumerate}
    \item Put $R=\Z/12\Z$. For which elements $f\in R$ is the localization $R_f= 0$?
    \item Put $R=\Z/12\Z$. For which elements $f\in R$ is the localization $R_f$ a field? (Recall that the zero ring is not a field.)
    \item The Jordan form of an operator $A:\C^6\to \C^6$ is
    \[
    \begin{pmatrix}
    0 & 1 & 0 & 0 & 0 & 0\\
    0 & 0 & 1 & 0 & 0 & 0\\
    0 & 0 & 0 & 1 & 0 & 0\\
    0 & 0 & 0 & 0 & 0 & 0\\
    0 & 0 & 0 & 0 & 1 & 1\\
    0 & 0 & 0 & 0 & 0 & 1
    \end{pmatrix}
    \]
    What is the Jordan form of the operator $A^3-A$?
\end{enumerate}
\end{question}
\begin{solution}
\begin{enumerate}
    \item $f = 0,6$, since these are the only nilpotent elements of $R$.
    \item We can answer this question by thinking about the prime ideals of $R_f$. We know $R$ has two prime ideals: $(2)$ and $(3)$. In order for $R_f$ to be a field, it needs to have exactly one prime ideal, which must be the zero ideal. We know the primes in the localization will be the prime ideals of $R$ which do not contain a power of $f$. A prime ideal contains a power of $f$ iff it contains $f$, so in order to eliminate one of the prime ideals we need to choose either $f \in (3)$ or $f \in (2)$. We already determined in part (a) that localizing at $f = 6$ will give us the zero ring, which is not a field. Localizing at an element of $(3)$ which is not in $(6)$ (so $f = 3$ or $f = 9$) will leave us with a ring with three ideals ($(1)$, $(2)$, and $(4) = (0)$), which is not a field. So, our other option is localizing at any element of $(2)$ which is not in $(6)$ ($f = 2,4,8,$ or $10$), and all of these do work.
    \item It is possible to do this problem computationally, but working with $6 \times 6$ matrices seems like a lot of work, so let's try to do this using other properties of JNF. From the Jordan form of $A$, we can see that its minimal polynomial is $f(x) = x^4(x-1)^2$. Thus, the minimal polynomial of $A^3-A$ is $x^4$, because in the ring $\C[x]/f(x)$, the polynomial $x^3-x = x(x-1)(x+1)$ is nilpotent of nilpotence degree 4 (i.e.\ $(x^3-x)^4 = 0$ but $(x^3-x)^k \neq 0$ for $k < 4$). This tells us that the only eigenvalue of $A^3-A$ is 0, and the size of the largest Jordan block is 4. This leaves us with only two possible Jordan forms for $A^3-A$: it has Jordan blocks of sizes $(4,2)$ or of sizes $(4,1,1)$.
    
    We can tell these two possibilities apart by checking the dimension of the 0-eigenspace (i.e. the kernel) of $A^3-A$: if it is 2-dimensional then we are in the former case, and if it is 3-dimensional we are in the latter case. The kernel of $A^3-A$ will be the sum of the kernels of $A$, $A-I$, and $A+I$ (this is true in general for polynomials in terms of matrices). That is, $\ker (A^3-A)$ will be the sum of the 0-, 1-, and $\shortminus 1$-eigenspaces of $A$. We can tell from the Jordan form of $A$ that the 0-eigenspace is 1-dimensional, as is the 1-eigenspace, and there are no eigenvectors of eigenvalue $\shortminus1$, so $\dim \ker (A^3-A) = 2$, and we are in the case where we have Jordan blocks of size $(4,2)$. Thus, the Jordan form of $A^3-A$ is
    \[
    \begin{pmatrix}
    0 & 1 & 0 & 0 & 0 & 0\\
    0 & 0 & 1 & 0 & 0 & 0\\
    0 & 0 & 0 & 1 & 0 & 0\\
    0 & 0 & 0 & 0 & 0 & 0\\
    0 & 0 & 0 & 0 & 0 & 1\\
    0 & 0 & 0 & 0 & 0 & 0
    \end{pmatrix}
    \]
\end{enumerate}
\end{solution}

%non-Galois extension, ideals of M_2(\Z), failure of tensor product to be exact, primes in a localization
% TENS, LOC
\begin{question}[subtitle=January 2021 Problem 1,ID=Jan21_1]
On this problem, no justification is required; only the answer will be graded.
\begin{enumerate}
    \item Give an example of a finite extension of $\Q$ which is not Galois.
    \item Give an example of two-sided ideal of the ring of $2 \times 2$ matrices over $\Z$ which is neither 0 nor the whole ring.
    \item Give an example of an injective map $M \to N$ of $k[x]$-modules ($k$ a field) and another $k[x]$-module $R$ such that $M\tensor R \to N \tensor R$ is not injective.
    \item Give an example of a commutative ring $R$ and prime ideals $p,q,q'$ in $R$ such that the extension $\widebar{q}$ of $q$ to the localization $R_p$ is prime, while the extension $\overline{q'}$ of $q'$ to $R_p$ is not prime.
\end{enumerate}
\end{question}
\begin{solution}
\begin{enumerate}
    \item $\Q(\sqrt[4]{2})$
    \item $M_2(2\Z)$
    \item Consider $k[x] \tonamed{\cdot x} k[x]$ and $R = k[x]/(x)$. Then $k[x]\tensor_{k[x]} R\cong R \tonamed{\cdot x} k[x]\tensor_{k[x]} R \cong R$ is not injective since it is the zero map.
    \item Let $R = \Z$, $p = (2)$, $q=(0)$, and $q'=(3)$.
\end{enumerate}
\end{solution}

%Field theory (any quadratic extension is given by adjoining a square root), ring theory (definition of an algebra/algebra hom), linear algebra (minimal polynomial, eigenvalues and diagonalizability/upper-triangularizability)
\begin{question}[subtitle = {January 2021 Problem 4, edited},ID=Jan21_4]
Let $F/\Q$ be a quadratic extension, and let $f$ and $g$ be two maps of $\Q$-algebras from $F$ to $M_2(\Q)$. (\textbf{NB} in this problem we take algebra homomorphisms to be \emph{unital}, meaning that $f(1) = g(1) = I$.)
\begin{enumerate}
    \item Prove that there exists a matrix $X \in \GL_2(\Q)$ such that for all $\alpha \in F$, $X f(\alpha) X\inv = g(\alpha)$.
    \item Prove that there does not exist a choice of $f$ whose image is contained in the set of upper triangular matrices.
\end{enumerate}
\end{question}
\begin{solution}
\begin{enumerate}
    \item For any $\Q$-algebra homomorphism $f\colon F \to M_2(\Q)$, we must have that $f(q) = qI$ for all $q \in \Q$. So, if $F = \Q(\sqrt{n})$, we need only show that there is an invertible matrix $X$ so that $X f(\sqrt{n}) X\inv = g(\sqrt{n})$. Let $A = f(\sqrt{n})$, then $A$ must satisfy $A^2 - nI = 0$, so the minimal polynomial of $A$ divides $x^2-n$. Since $x^2-n$ is irreducible over $\Q$, in fact it must be the minimal polynomial of $A$, so the eigenvalues of $A$ are $\pm \sqrt{n}$. Thus, over $\C$ we can find an invertible matrix $Y$ so that
    \[YAY\inv = \twomatrix{\sqrt{n}}{0}{0}{\shortminus \sqrt{n}}.\]
    The same is true for $B = g(\sqrt{n})$, so this shows that over $\C$ we can find an $X$ so that $XAX\inv = B$. Then $XA - BX = 0$. Notice that $T \mapsto TA - BT$ is a linear transformation $M_2(\Q) \to M_2(\Q)$, and we have just shown that it has nontrivial kernel as a linear transformation $M_2(\C) \to M_2(\C)$. Therefore, it must have nontrivial kernel as a linear transformation $M_2(\Q) \to M_2(\Q)$, and so we can take $X \in \GL_2(\Q)$.

    \item As noted above, $A = f(\sqrt{n})$ has eigenvalues $\pm \sqrt{n}$. If $A$ were upper triangular, its diagonal entries would have to be its eigenvalues, but $\pm \sqrt{n} \not\in \Q$. 
\end{enumerate}
\end{solution}

%LGFF. non-diagonalizable matrix, normal subgroup of matrix group, one-sided ideal in M_2(\C), flat module that is not free, definition of splitting of exact sequence
% FLAT
\begin{question}[subtitle=August 2020 Problem 1,ID=Aug20_1]
On this problem, no justification is required; only the answer will be graded.
\begin{enumerate}
    \item Give an example of a non-diagonalizable matrix in $\GL_3(\F_p)$ (for some $p$).
    \item Give an example of a nontrivial proper normal subgroup of $\GL_3(\F_p)$ (for some $p$).
    \item Give an example of a one-sided ideal (either left or right) in $M_2(\C)$ which is not a two-sided ideal.
    \item Give an example of a $\Z$-module which is flat but not free.
    \item Give an example of an exact sequence of abelian groups which is not split.
\end{enumerate}
\end{question}
\begin{solution}
\begin{enumerate}
    \item Let $p=2$ and consider $M = \begin{pmatrix} 1 & 1 & 0\\ 0 & 1 & 1 \\ 0 & 0 & 1\end{pmatrix}$.
    \item Let $p=3$. Then $\SL_3(\F_3)$ is a proper normal subgroup of $\GL_3(\F_3)$.
    \item Let $I$ be the ideal of matrices of the form $\begin{pmatrix} a & 0\\ b & 0\end{pmatrix}$. Then $I$ is a left ideal but not a right ideal.
    \item $\Q$
    \item $0 \to \Z \to \Z \to \Z/2\Z \to 0$
\end{enumerate}
\end{solution}

%Jan '20 is about the ring theory of \Z[S_3]

%
\begin{question}[subtitle = August 2019 Problem 1,ID=Aug19_1]
On this problem, only the answers will be graded.
\begin{enumerate}
    \item Let $R = M_2(\Z)$ the ring of $2 \times 2$ matrices with entries in $\Z$. Give an example of a left ideal that is not a right ideal.
    \item Give an example in $R = M_2(\Z)$ of a proper, nonzero, two-sided ideal.
    \item $\Z/100\Z \tensor_\Z \Z/35\Z$ is a finite abelian group. Compute its order.
\end{enumerate}
\end{question}
\begin{solution}
\begin{enumerate}
    \item Let $I$ be the ideal of matrices of the form $\begin{pmatrix} a & 0\\ b & 0\end{pmatrix}$. Then $I$ is a left ideal but not a right ideal.
    \item The set of matrices with even entries is a proper, nonzero, two-sided ideal of $M_2(\Z)$.
    \item We have that $\Z/100\Z \tensor_\Z \Z/35\Z \cong \Z/\gcd(100,35)\Z = \Z/5\Z$, so the tensor product has order 5.
\end{enumerate}
\end{solution}

%Field extensions + faithful actions
% TENS
\begin{question}[subtitle = {August 2019 Problem 3, edited},ID=Aug19_3]
Let $L/K$ be a Galois extension of fields with Galois group $G$. Let $S$ be the subset of roots of unity in $L$ other than 1; that is, $S$ is the set of elements $x\in L$ such that $x \neq 1$ but $x^k = 1$ for some positive integer $k$.
\begin{enumerate}
    \item Show that the action of $G$ on $L$ restricts to an action of $G$ on $S$ (i.e., the action of $G$ preserves $S$.)
    \item Let $L=\Q(i)$ and $K=\Q$. Show in this case that $G$ acts faithfully on $S$.
    \item Give an example of $L/K$ such that the action of $G$ on $S$ is not faithful.
    \item Suppose that the action of $G$ on $S$ is transitive. Prove that $\mathrm{char}(K)=2$.
\end{enumerate}
\end{question}
\begin{solution}
\begin{enumerate}
    \item The $n$th roots of unity in $L$ are precisely the elements in $L$ that are roots of the equations $t^n-1$, and the roots of unity different from 1 are the roots of $(t^n-1)/(t-1)$. Since $(t^n-1)/(t-1) \in K[t]$, $G$ must permute the roots of this polynomial. Thus, $G$ also acts on $S$.
    \item In this case, $S = \set{\shortminus 1, \pm i}$, and $G$ is cyclic of order 2 generated by complex conjugation $\sigma$. So the action of $G$ on $S$ is given by $\id$ acting trivially and $\sigma$ acting by $i \mapsto -i$. So this action is faithful since only the identity element in $G$ acts trivially on $S$.
    \item Let $L = \Q(\sqrt{2})$ and $K = \Q$. Then $S = \set{\shortminus 1}$ since $\pm 1$ are the only roots of unity in $\R$, and $L \subseteq \R$. As in the previous example, $G$ is a cyclic group of order 2. But both elements of $G$ fix $S$ since $S \subseteq \Q$.
    \item Imagine $K$ has characteristic not equal to 2. Then $\shortminus 1$ and $1$ are distinct elements of $S \cap K$. But every element of $G$ will fix $\shortminus 1$ since it is in $K$, so the action could not be transitive.
\end{enumerate}
\end{solution}

%The one w/ the Smith Normal Form question
% PROJ kinda
\begin{question}[subtitle=January 2019 Problem 1,ID=Jan19_1]
On this problem, only the answers will be graded.
\begin{enumerate}
    \item Give an example of a simple ring which is not a commutative ring.
    \item Give an example a ring $R$, $R$-modules $M_1, M_2, M_3$, and $N$ and a short exact sequence $0\to M_1\to M_2 \to M_3 \to 0$  such that the corresponding sequence \newline$0\leftarrow \Hom(M_1, N)\leftarrow \Hom(M_2, N)\leftarrow \Hom(M_3, N)\leftarrow 0$ is not exact.
    \item Let $\Phi$ be the map $\Z^3\to \Z^2$ given by the matrix $\begin{bmatrix} 2 & 3 & 5\\ 12 & \shortminus 4 & 8 \end{bmatrix}$. Let $M$ be the cokernel of $\Phi$, which we recall is $\Z^2/\im(\Phi)$. Compute the annihilator of $M$.
\end{enumerate}
\end{question}
\begin{solution}
\begin{enumerate}
    \item Let $R$ be the ring of $2 \times 2$ matrices over $\F_2$. Then $R$ is not commutative, and $R$ is simple since the only nonzero two-sided ideal is $R$.
    \item Let $R = \Z$ and consider $0 \to \Z \to \Z \to \Z/2\Z \to 0$. Then applying $\Hom(-,\Z/2\Z)$ gives
    \[ 0 \to \Hom(\Z/2\Z,\Z/2\Z) \to \Hom(\Z, \Z/2\Z)\cong \Z/2\Z \to \Hom(\Z,\Z/2\Z)\cong \Z/2\Z \to 0\]
    which is not exact on the right since the map is multiplication by $2=0$, so it is not surjective.
    \item I guess this space will be where I write about Smith Normal Form. Notice that computing the cokernel of $\Phi$ would be easy if we had a matrix of the form
    \[\begin{bmatrix} a & 0 & 0\\ 0 & b & 0 \end{bmatrix},\]
    as then the cokernel would clearly be $\Z/a\Z \oplus \Z/b\Z$. The goal of Smith Normal Form is to rewrite a given matrix in a new pair of bases so that it looks like this.

    Computing SNF is very similar to row-column reduction over a field, with two notable caveats:
    \begin{enumerate}[(i)]
        \item You cannot multiply or divide rows/columns by scalars (other than $\pm 1$). Instead, your only available techniques are to swap rows/colums, and to add multiples of one row/column to another. You can use this to replace an element with the gcd of all the elements in its row/column, and then clear those other elements as usual.
        \item Once you clear all the entries in the column below a pivot, you might have to ``ruin'' your work in order to clear the entries to the right in the row, and visa versa. However, by alternately clearing the entries below and to the right, eventually you will be able to clear both. (Luckily, this process doesn't actually come up in this particular problem.)
    \end{enumerate}
    We begin with the matrix from the problem. I notice that in the first row, the elements have gcd 1, so I'm going to subtract the first column from the second column, then swap the first two columns:
    \[\begin{bmatrix} 1 & 2 & 5\\ \shortminus 16 & 12 & 8 \end{bmatrix}.\]
    Now I can use the 1 in the first row to clear the $\shortminus 16$ below it, and the 2 and 5 to the right of it (the order that I do this does not matter, because it corresponds to multiplying by matrices on either side, and matrix multiplication is associative). Once I do, I arrive at this matrix:
    \[\begin{bmatrix} 1 & 0 & 0\\ 0 & 44 & 88 \end{bmatrix}.\]
    Luckily, the process ends right away, since I can use the 44 to clear the 88, so the SNF is:
    \[\begin{bmatrix} 1 & 0 & 0 \\ 0 & 44 & 0 \end{bmatrix}.\]
    Thus, the cokernel of $\Phi$ is isomorphic to $\Z/44\Z$, and the annihilator is 44.

    (Here's another true thing: the annihilator of the cokernel of $\Phi$ is equal to the gcd of the maximal subdeterminants of $\Phi$. In this case, the subdeterminants are $\shortminus 44$, $\shortminus 44$ and $44$, and the gcd of those numbers is easy to find. In general, this is a strictly more cumbersome algorithm than computing the Smith Normal Form directly, but it works nicely when you only have a small number of $2 \times 2$ subdeterminants.)
\end{enumerate}
\end{solution}

%This is actually a field theory/Galois theory question. It's not very hard, actually, and I like that it involves a little bit of thinking about splitting/ramifying/extending behaviors.
\begin{question}[subtitle = August 2018 Problem 5,ID=Aug18_5]
Let $A$ be a two-dimensional (unital) algebra over a field $F$. This means $A$ is an associative, but not necessarily commutative, ring with a unit that contains $F$ as a subring, such that the elements of $F$ commute with the elements of $A$ (that is, $F$ is in the center) and $A$ is two-dimensional as a vector space over $F$.
\begin{enumerate}
    \item Show that $A$ must in fact be commutative.
    \item Show that if $F$ is algebraically closed, then either $A \simeq F \times F$ or $A \simeq F[x]/(x^2)$.
    \item Suppose $F = \R$. List (with a proof) all the possibilities for $A$, up to isomorphism.
\end{enumerate}
\end{question}
\begin{solution}
\begin{enumerate}
    \item Since $A$ is 2-dimensional as an $F$-vector space, and contains $F$, we can choose a basis of the form $1,a$ for some $a \in A$. Then any element of $A$ can be written as $x+ya$ with $x,y \in F$. Since 1 and $a$ certainly commute with one another, in fact all the elements of $A$ commute with one another.
    \item Since $A$ is a commutative $F$-algebra generated as an algebra by the element $a$, there is a map $F[x] \to A$ sending $x \to a$, and the kernel of this map is an ideal of $F[x]$. Such an ideal is principal, since $F[x]$ is a PID, and its generator must be a degree 2 polynomial, since $A$ is 2-dimensional. As $F$ is algebraically closed, this polynomial factors into linear terms. If it factors into distinct linear terms, then we get $A \simeq F[x]/(x-a)(x-b) \simeq F[x]/(x-a) \times F[x]/(x-b)$ by the Chinese Remainder Theorem, so $A \simeq F \times F$. On the other hand, if it has a repeated root, then $A \simeq F[x]/((x-a)^2) \simeq F[x]/(x^2)$.
    \item In this case, in addition to the two possibilities in part (b), we could also have that the polynomial is irreducible over $\R$. In such a case, by the fundamental theorem of algebra we get that $A \simeq \C$.
\end{enumerate}
\end{solution}

%Aug '18 has three ring-ish questions

%Jan '18 has three questions about localizations

%Aug '17 has the three neat tensor product ones

%Jan '17 is all group theory

%It was pointed out to me that probably the intended solution is to use RCF. The argument I have below is essentially the same, but substituting some field theory facts in for knowledge of RCF.
\begin{question}[subtitle = August 2017 Problem 2,ID=Aug17_2]
Let $K$ be a field, and let $A$ be an $n\times n$ matrix over $K$. Suppose that $f(x) \in K[x]$ is an irreducible polynomial such that $f(A)=0$. Show that $\deg(f) \mid n$.
\end{question}
\begin{solution}
Let $m(x) \in K[x]$ be the minimal polynomial of $A$. Since $f(A)=0$, $m(x)$ must divide $f(x)$, but $f(x)$ is irreducible so $m(x)=f(x)$. Let $c(x)$ be the characteristic polynomial of $A$, which has degree $n$. Since $c(x)$ and $f(x)$ have the same roots and $f(x)$ is irreducible, it must be the case that $c(x) = f(x)^\ell$ for some $\ell$. Therefore, $\ell\cdot\deg(f) = \deg(c) = n$, so $\deg(f) \mid n$.
\end{solution}

%Prior to 2017, there were not many "no justification mix" kind of problems.

%LG
\begin{question}[subtitle=January 2016 Problem 5,ID=Jan16_5]
Consider the group $\GL_n(\Q)$ of invertible $n\times n$ matrices with rational coefficients. Suppose $G \subset \GL_n(\Q)$ is a finite subgroup. Prove that every prime factor $p$ of the order $\abs{G}$ satisfies $p \leq n+1$.
\end{question}
\begin{solution}
Suppose that $p \mid \abs{G}$, then by Cauchy's theorem there is an element $A \in G$ of order $p$, $A^p = I$. Notice that the minimal polynomial of $A$ divides $x^p - 1$, so the minimal polynomial of $A$ is either $x^p-1$ or $\Phi_p(x)$ (notice it cannot be $x-1$). Since the minimal polynomial divides the characteristic polynomial, the degree of the minimal polynomial is at most $n$. Thus, if the minimal polynomial is $x^p-1$ then we get $p \leq n$, and if the minimal polynomial is $\Phi_p(x)$ we get $p-1 \leq n$. Thus, in all cases $p \leq n+1$.

I was curious whether or not the bound on $p$ can actually be achieved, and the answer is it can. For example, the $2\times 2$ matrix
\[A = \twomatrix{0}{-1}{1}{-1}\]
has order 3.
\end{solution}